
The rapid evolution of enterprise web applications has driven a paradigm shift from monolithic architectures to cloud-native \textit{microservices} solutions. This transition is motivated by the need for greater scalability, resilience, and agility in the software development life cycle. This thesis presents the design, development and deployment of "\textit{Smart Travel}", a comprehensive e-commerce platform dedicated to the sale of travel packages. The project serves as a proof of concept (\gls{PoC}) developed within an industrial internship context, aiming to analyze the complexities and benefits of a modern and distributed software architecture.

The primary objective of this work is not merely the implementation of functional requirements (such as catalog management and order processing), but rather the critical evaluation of competing technologies and the automation of the release process. A central element of this research is the comparative analysis between two leading Java ecosystems: the established \textit{Spring Boot} framework and the cloud-native \textit{Quarkus} framework. Specifically, the project adopts a fully reactive approach for both of them, utilizing \textit{Spring WebFlux} and \textit{Quarkus Reactive} capabilities.

From an architectural perspective, the Smart Travel system is designed to ensure loose coupling and high availability. The backend adopts a microservices pattern supported by \textit{RabbitMQ}, a message broker that facilitates asynchronous \textit{event-driven} communication among services. Data persistence is managed through \textit{MongoDB}, a \gls{NoSQL} database chosen for its schema flexibility, which is particularly advantageous for handling diverse and evolving travel package structures compared to rigid relational models. On the client side, the user interface is delivered via a \textit{Single Page Application (\gls{SPA})} developed with \textit{Angular} and leveraging \textit{\gls{RxJS}} for reactive state management, ensuring a responsive and fluid user experience.

A significant portion of this thesis is dedicated to the "industrialization" of the software, bridging the gap between development and operations (\textit{DevOps}). The project moves beyond local execution by implementing a robust \textit{Continuous Integration and Continuous Delivery (\gls{CI/CD})} pipeline using \textit{Jenkins}. This pipeline automates the entire delivery chain, including code checkout, compilation, container image creation and pushing to a private registry. The final deployment target is \textit{OKD} (the community distribution of \textit{Kubernetes} that powers \textit{Red Hat OpenShift}). The thesis details the orchestration of containerized applications within the OKD cluster, managing critical aspects such as service discovery, route exposure and configuration management via \texttt{ConfigMaps} and \texttt{Secrets}.

In conclusion, the Smart Travel project demonstrates the viability and robustness of a heterogeneous, fully reactive microservices architecture. It successfully integrates different modern technologies into a cohesive ecosystem, proving that the adoption of cloud-native principles, while introducing complexity in orchestration and consistency management, offers superior long-term benefits in terms of maintainability and scalability. The work provides a practical blueprint for modernizing legacy Java applications, offering concrete evidence of how the synergy between Reactive frameworks, Event-Driven patterns, and automated DevOps practices defines the state of the art in contemporary software engineering.